\documentclass[11pt]{article}

\usepackage[a4paper,margin=1in]{geometry}
\usepackage[T1]{fontenc}
\usepackage[utf8]{inputenc}
\usepackage{lmodern}
\usepackage{booktabs}
\usepackage{longtable}
\usepackage{hyperref}
\usepackage{enumitem}

\title{Clinician-Readiness Dataset Process Report\\v0.3 Final Line}
\author{FSD Mental Health Safety Benchmark Team}
\date{17 February 2026}

\begin{document}
\maketitle

\section{Purpose}
This report documents the final process that led to the current clinician-readiness dataset line in this repository. The intent is to describe why the changes were made and what changed at study level, not to justify individual commits.

The process is framed as three phases:
\begin{enumerate}[leftmargin=1.5em]
  \item small-scope benchmark run,
  \item scaling to publication-grade sample sizes,
  \item clinician-readiness hardening and reproducibility gates.
\end{enumerate}

\section{Claim Boundary and Positioning}
OpenR1-Psy is used as the upstream source corpus and is treated as clinically grounded according to the published construction process (DSM/ICD-guided reasoning plus automated filtering). This project does \textbf{not} claim full clinician adjudication of all labels/plans in the current release line.

The v0.3 work therefore focuses on:
\begin{itemize}[leftmargin=1.5em]
  \item reproducibility (frozen snapshots, local hash manifests, release discipline),
  \item benchmark integrity (label/mapping/schema consistency),
  \item clinician usability (join-free exports with explicit review metadata),
  \item fail-closed validation behaviour.
\end{itemize}

\section{Process Timeline}
\subsection{Phase 0: Small-Scope Baseline}
The initial benchmark documentation records a small run configuration (Study A around 300 prompts, Study B around 276 single-turn plus 10 multi-turn, Study C around 30 cases). This phase established metric wiring and runner behaviour before full scaling.

\subsection{Phase 1: Scaling}
Scaling moved the benchmark to larger operational sample sizes used by the current branch:
\begin{itemize}[leftmargin=1.5em]
  \item Study A split: 2000 samples,
  \item Study B split: 2000 single-turn samples,
  \item Study B multi-turn split: 120 cases (20-turn conversations),
  \item Study C split: 100 cases.
\end{itemize}

This stage was about statistical adequacy, persona coverage, and stress-testing conditions, not clinician package readiness.

\subsection{Phase 2: Clinician-Readiness}
The clinician-readiness line then introduced auditability and packaging controls. Commit sequence on this branch shows a clean evolution:
\begin{itemize}[leftmargin=1.5em]
  \item pre-update freeze capture,
  \item gold artefact rerun and diff capture,
  \item v0.3 frozen snapshot with declared artefacts,
  \item clinician package builder and stage gates,
  \item final tightening pass (schema, routing, preflight robustness).
\end{itemize}

\section{Study-Wise Updates}
\subsection{Study A: Label and Metadata Integrity}
\textbf{What changed}
\begin{itemize}[leftmargin=1.5em]
  \item Frozen v0.2 to v0.3 comparison shows 560/2000 label value differences in \texttt{gold\_diagnosis\_labels.json}.
  \item Tier-2 correction set applied for 8 audited IDs (non-clinical and substance-primary cases).
  \item Mapping coherence repaired on the legacy 300-row mapping surface (176 entries synced).
  \item Metadata remained sparse by design (27 explicit IDs), with 10 enriched entries and explicit review/certainty fields.
  \item Export canonicalisation added through \texttt{label\_canonical\_map.json} (17 canonical labels, 6 aliases).
\end{itemize}

\textbf{Why it was needed}
\begin{itemize}[leftmargin=1.5em]
  \item Remove label provenance conflicts that break audit defensibility.
  \item Prevent format-only label drift in clinician-facing exports.
  \item Keep uncertainty explicit for high-risk review rows.
\end{itemize}

\subsection{Study B: Schema Safety and Pressure Construct Reliability}
\textbf{What changed}
\begin{itemize}[leftmargin=1.5em]
  \item Runtime schema validation was hardened to return structured failure on malformed payloads (instead of raising).
  \item Multi-turn pressure set remained balanced (40/40/40 by style and schedule).
  \item Deterministic \texttt{turns\_text} export was added for spreadsheet review.
  \item Stage-2 validators were attached to package gates:
  \begin{itemize}[leftmargin=1.5em]
    \item multi-turn uniqueness (no duplicate full-turn signatures),
    \item construct validation (pass rate 97.5\% on 2000 rows).
  \end{itemize}
\end{itemize}

\textbf{Why it was needed}
\begin{itemize}[leftmargin=1.5em]
  \item Prevent runtime/schema edge cases from surfacing only after generation.
  \item Ensure pressure tests remain credible and non-template-like.
  \item Make clinician sheet review operational without JSON parsing.
\end{itemize}

\subsection{Study C: Plan Quality and Entity Anchoring Defensibility}
\textbf{What changed}
\begin{itemize}[leftmargin=1.5em]
  \item Frozen v0.2 to v0.3 comparison shows 92/100 target-plan entries changed.
  \item Entity evidence map added (100 case entries, 2 semantic synonym groups).
  \item Anchoring validator logic hardened:
  \begin{itemize}[leftmargin=1.5em]
    \item case evidence must be non-empty and present in \texttt{patient\_summary},
    \item synonym miss must still evaluate case-evidence fallback.
  \end{itemize}
  \item Study C exports include explicit \texttt{persona\_id} for duplicated-persona interpretation.
\end{itemize}

\textbf{Why it was needed}
\begin{itemize}[leftmargin=1.5em]
  \item Stop false-positive ``anchored'' outcomes from weak evidence handling.
  \item Keep plan/entity artefacts inspectable at case level.
  \item Avoid ambiguity in duplicated-persona review scenarios.
\end{itemize}

\section{Reproducibility and Release Discipline}
\subsection{Snapshot and Release Separation}
Two contracts are kept separate:
\begin{itemize}[leftmargin=1.5em]
  \item \textbf{Frozen snapshot source-of-truth}: \texttt{data/frozen\_splits/v0.3\_postclinician\_audit/}
  \item \textbf{Consumer release source-of-truth}: \texttt{data/releases/clinician\_readiness\_v0.3\_2026-02-16/}
\end{itemize}

\subsection{Deterministic Package Outputs}
Current clinician package row counts are fixed by builder gates:
\begin{itemize}[leftmargin=1.5em]
  \item Study A review: 2000 rows,
  \item Study B single-turn review: 2000 rows,
  \item Study B multi-turn review: 120 rows,
  \item Study C review: 100 rows,
  \item Safety-priority review: 32 rows.
\end{itemize}

\subsection{Preflight and Gate Status}
The v0.3 preflight report records all checks passing, including:
\begin{itemize}[leftmargin=1.5em]
  \item runtime schema tests,
  \item stage-2 gate script,
  \item package build and manifest checks,
  \item frozen snapshot manifest checks,
  \item release manifest verify checks.
\end{itemize}

\section{What Changed vs What Did Not}
\subsection{Changed}
\begin{itemize}[leftmargin=1.5em]
  \item Reference labels/plans/metadata consistency and review policy overlays.
  \item Validator strictness, packaging semantics, and reproducibility controls.
  \item Clinician-facing ergonomics (explicit columns, deterministic fields, routeable safety sheet).
\end{itemize}

\subsection{Not Changed}
\begin{itemize}[leftmargin=1.5em]
  \item No raw split regeneration in the final tightening pass.
  \item No schema rewrite of upstream raw OpenR1-Psy payload files.
  \item No claim of full clinician adjudication for the entire corpus.
\end{itemize}

\section{Decision Rationale for Keeping v0.3}
Keeping the v0.3 line is the defensible stop point for this cycle:
\begin{itemize}[leftmargin=1.5em]
  \item it preserves comparability with upstream OpenR1-Psy content,
  \item it materially improves benchmark reliability and reproducibility,
  \item it avoids overclaiming clinician validation,
  \item it provides a clean base for optional future clinician subset adjudication.
\end{itemize}

\section{Future Work (Optional)}
\begin{itemize}[leftmargin=1.5em]
  \item stratified clinician review subset for reference-quality calibration,
  \item expansion from sparse to dense Study A metadata if operationally required,
  \item further Study C entity-count naturalisation beyond the v0.3 normalised baseline.
\end{itemize}

\section{Primary Artefact References}
\begin{itemize}[leftmargin=1.5em]
  \item \texttt{benchmark/runtime/data/frozen\_splits/v0.2\_preclinician/manifest.json}
  \item \texttt{benchmark/runtime/data/frozen\_splits/v0.3\_postclinician\_audit/manifest.json}
  \item \texttt{benchmark/runtime/data/frozen\_splits/v0.3\_postclinician\_audit/audit\_fix\_diff\_summary.json}
  \item \texttt{benchmark/runtime/docs/reports/clinician\_package/v0.3/sendoff\_preflight\_report.json}
  \item \texttt{benchmark/runtime/docs/reports/clinician\_package/v0.3/stage2\_gate\_report.json}
  \item \texttt{benchmark/runtime/data/releases/LATEST.md}
\end{itemize}

\end{document}
